\movedsub{Direction is dataset-dependent}
\label{sec:direction}

Where a difference exists it does not point the same way. Correlating NDCG@5 with mean faithfulness across the four systems, in all twenty-four dataset-generator-evaluator combinations (Appendix~\ref{app:spearman}), splits the two datasets. Under AlignScore the split is exact: all four HotpotQA coefficients are positive, reaching +1.00 in three of them, and all four Natural Questions coefficients are zero or negative. Under NLI-max it holds in seven of the eight. Every HotpotQA coefficient is positive and three of the four Natural Questions coefficients are negative, reaching -1.00 for Claude. Mistral-7B on Natural Questions is the exception, at +0.80. Under claim-min the split fails in four of eight, and all four failures fall on the two open-weight generators. Better retrieval accompanies more faithful answers on the multi-hop dataset and not on the single-hop one. Composing an answer out of several retrieved pieces plausibly makes grounding sensitive to how much of the evidence is present, where a single-fact answer is not. We do not test that. The experiment varies embedding model rather than manipulating retrieval quality, so none of this is an effect of NDCG@5 itself.

Claim-min breaks the split. Four of its eight combinations go the wrong way, and all four are on the two open-weight generators. Qwen inverts the split on both datasets, at +0.63 on Natural Questions and -0.40 on HotpotQA. Mistral-7B is positive on Natural Questions, at +0.80, and flat on HotpotQA. Drop claim-min and fifteen of the remaining sixteen combinations follow the split; drop the two open-weight generators instead and all twelve do. The third evaluator and the fourth generator are what show the direction is not a property of the dataset alone.
